% --- Archon physics formalization source begin ---
% archon:physics
% archon:covers PhyXMiniProblems/problem_phyx_mini_0244.lean
% archon:source-report reports/phyx_mini/problem_phyx_mini_0244.source.json

\chapter{Physics problem phyx\_mini\_0244}
\label{ch:PhyXMiniProblems_problem_phyx_mini_0244}

\paragraph{Problem source.}
\#\# Physical scenario

If a loop of chain is spun at high speed, it can roll along the ground like a circular hoop without collapsing. Consider a chain of uniform linear mass density \textbackslash{}( \textbackslash{}mu \textbackslash{}) whose center of mass travels to the right at a high speed \textbackslash{}( v\_0 \textbackslash{}) as shown in figure.Assume the weight of an individual link is negligible compared to the tension.

\#\# Question

Determine the tension in the chain in terms of \textbackslash{}( \textbackslash{}mu \textbackslash{}) and \textbackslash{}( v\_0 \textbackslash{}).

\#\# Answer choices

- A: A: \$T = \textbackslash{}alpha \textbackslash{}sqrt\{v\_0\}\$

- B: B: \$T = \textbackslash{}alpha v\_0\textasciicircum{}3\$

- C: C: \$T = \textbackslash{}alpha v\_0\$

- D: D: \$T = \textbackslash{}alpha v\_0\textasciicircum{}2\$

\#\# Figure caption (auxiliary; use the image as primary evidence)

The image depicts a chain arranged in a circular shape on a horizontal surface. There's a red arrow labeled \textbackslash{}( \textbackslash{}vec\{v\}\_0 \textbackslash{}) pointing to the right, indicating initial velocity. On the surface, there's a small protrusion labeled "Bump."

\#\# Dataset metadata

Wave Properties; Physical Model Grounding Reasoning, Multi-Formula Reasoning

\paragraph{Recorded answer/context.}
D: D: \$T = \textbackslash{}alpha v\_0\textasciicircum{}2\$

\paragraph{Figure/image path.}
/root/proposal\_for\_physic/science-mango/phyx\_mini\_run/phyx\_data/test\_image/244.png

\paragraph{Formalization target.}
create a compiling Lean file with sorry bodies at `PhyXMiniProblems/problem\_phyx\_mini\_0244.lean`. The Lean declarations must preserve the physical quantities, dimensions or dimensional roles, figure labels, governing-law hypotheses, and final relation expressed by this problem.
Use Mathlib/Physlib names found through LeanExplore where available. If a physics API is missing, introduce faithful local abstractions rather than scalar placeholder aliases.

\begin{theorem}[Physics formalization target]
\label{thm:physics:phyx_mini_0244:target}
\lean{PhyXMiniProblems.ProblemPhyXMini0244.problem_phyx_mini_0244}
\uses{def:physics:phyx-mini-0244:phyxminiproblems-problemphyxmini0244-linearmassdensityreadout, def:physics:phyx-mini-0244:phyxminiproblems-problemphyxmini0244-speedreadout, def:physics:phyx-mini-0244:phyxminiproblems-problemphyxmini0244-forcereadout, def:physics:phyx-mini-0244:phyxminiproblems-problemphyxmini0244-rollingchainsetup, def:physics:phyx-mini-0244:phyxminiproblems-problemphyxmini0244-matchesproblemstatement, def:physics:phyx-mini-0244:phyxminiproblems-problemphyxmini0244-matchessuppliedfigure, def:physics:phyx-mini-0244:phyxminiproblems-problemphyxmini0244-hasphysicalrollingparameters, def:physics:phyx-mini-0244:phyxminiproblems-problemphyxmini0244-satisfiesrollingchainlaws}
The assigned autoformalize agent should translate this physics problem into Lean declarations in the covered file, with theorem and lemma proof bodies written as `by sorry`.
\end{theorem}
\begin{proof}
This is an autoformalization task, not a proof task. Produce statements that can later be proved by the physics prover without weakening the physical model.
\end{proof}

% --- Archon declaration topology begin ---
\section{Declaration topology}
\begin{definition}[Chain Radius]
  \label{def:physics:phyx-mini-0244:phyxminiproblems-problemphyxmini0244-chainradius}
  \lean{PhyXMiniProblems.ProblemPhyXMini0244.ChainRadius}
  The physical radius of the circular chain loop
\end{definition}
\begin{proof}
  This is the declared model definition of Chain Radius.
\end{proof}
\begin{definition}[Linear Mass Density]
  \label{def:physics:phyx-mini-0244:phyxminiproblems-problemphyxmini0244-linearmassdensity}
  \lean{PhyXMiniProblems.ProblemPhyXMini0244.LinearMassDensity}
  Uniform chain mass per unit arc length, with dimension mass / length
\end{definition}
\begin{proof}
  This is the declared model definition of Linear Mass Density.
\end{proof}
\begin{definition}[Angular Speed]
  \label{def:physics:phyx-mini-0244:phyxminiproblems-problemphyxmini0244-angularspeed}
  \lean{PhyXMiniProblems.ProblemPhyXMini0244.AngularSpeed}
  Angular speed of the rolling loop, with dimension 1 / time
\end{definition}
\begin{proof}
  This is the declared model definition of Angular Speed.
\end{proof}
\begin{definition}[Force Magnitude]
  \label{def:physics:phyx-mini-0244:phyxminiproblems-problemphyxmini0244-forcemagnitude}
  \lean{PhyXMiniProblems.ProblemPhyXMini0244.ForceMagnitude}
  A force magnitude, used for both tension and individual-link weight
\end{definition}
\begin{proof}
  This is the declared model definition of Force Magnitude.
\end{proof}
\begin{definition}[radius Readout]
  \label{def:physics:phyx-mini-0244:phyxminiproblems-problemphyxmini0244-radiusreadout}
  \lean{PhyXMiniProblems.ProblemPhyXMini0244.radiusReadout}
  \uses{def:physics:phyx-mini-0244:phyxminiproblems-problemphyxmini0244-chainradius}
  Scalar radius readout in a selected coherent unit system
\end{definition}
\begin{proof}
  This is the declared model definition of radius Readout.
\end{proof}
\begin{definition}[linear Mass Density Readout]
  \label{def:physics:phyx-mini-0244:phyxminiproblems-problemphyxmini0244-linearmassdensityreadout}
  \lean{PhyXMiniProblems.ProblemPhyXMini0244.linearMassDensityReadout}
  \uses{def:physics:phyx-mini-0244:phyxminiproblems-problemphyxmini0244-linearmassdensity}
  Scalar linear-mass-density readout in a selected coherent unit system
\end{definition}
\begin{proof}
  This is the declared model definition of linear Mass Density Readout.
\end{proof}
\begin{definition}[speed Readout]
  \label{def:physics:phyx-mini-0244:phyxminiproblems-problemphyxmini0244-speedreadout}
  \lean{PhyXMiniProblems.ProblemPhyXMini0244.speedReadout}
  Scalar speed readout in a selected coherent unit system
\end{definition}
\begin{proof}
  This is the declared model definition of speed Readout.
\end{proof}
\begin{definition}[angular Speed Readout]
  \label{def:physics:phyx-mini-0244:phyxminiproblems-problemphyxmini0244-angularspeedreadout}
  \lean{PhyXMiniProblems.ProblemPhyXMini0244.angularSpeedReadout}
  \uses{def:physics:phyx-mini-0244:phyxminiproblems-problemphyxmini0244-angularspeed}
  Scalar angular-speed readout in a selected coherent unit system
\end{definition}
\begin{proof}
  This is the declared model definition of angular Speed Readout.
\end{proof}
\begin{definition}[force Readout]
  \label{def:physics:phyx-mini-0244:phyxminiproblems-problemphyxmini0244-forcereadout}
  \lean{PhyXMiniProblems.ProblemPhyXMini0244.forceReadout}
  \uses{def:physics:phyx-mini-0244:phyxminiproblems-problemphyxmini0244-forcemagnitude}
  Scalar force readout in a selected coherent unit system
\end{definition}
\begin{proof}
  This is the declared model definition of force Readout.
\end{proof}
\begin{definition}[Horizontal Direction]
  \label{def:physics:phyx-mini-0244:phyxminiproblems-problemphyxmini0244-horizontaldirection}
  \lean{PhyXMiniProblems.ProblemPhyXMini0244.HorizontalDirection}
  Horizontal direction of the velocity arrow drawn in the figure
\end{definition}
\begin{proof}
  This is the declared model definition of Horizontal Direction.
\end{proof}
\begin{definition}[Loop Profile]
  \label{def:physics:phyx-mini-0244:phyxminiproblems-problemphyxmini0244-loopprofile}
  \lean{PhyXMiniProblems.ProblemPhyXMini0244.LoopProfile}
  Shape of the chain's visible centerline in the primary image
\end{definition}
\begin{proof}
  This is the declared model definition of Loop Profile.
\end{proof}
\begin{definition}[Supporting Surface Orientation]
  \label{def:physics:phyx-mini-0244:phyxminiproblems-problemphyxmini0244-supportingsurfaceorientation}
  \lean{PhyXMiniProblems.ProblemPhyXMini0244.SupportingSurfaceOrientation}
  Orientation of the supporting surface drawn below the chain
\end{definition}
\begin{proof}
  This is the declared model definition of Supporting Surface Orientation.
\end{proof}
\begin{definition}[Ground Feature Label]
  \label{def:physics:phyx-mini-0244:phyxminiproblems-problemphyxmini0244-groundfeaturelabel}
  \lean{PhyXMiniProblems.ProblemPhyXMini0244.GroundFeatureLabel}
  Text label attached to the small protrusion on the ground
\end{definition}
\begin{proof}
  This is the declared model definition of Ground Feature Label.
\end{proof}
\begin{definition}[Chain Motion Regime]
  \label{def:physics:phyx-mini-0244:phyxminiproblems-problemphyxmini0244-chainmotionregime}
  \lean{PhyXMiniProblems.ProblemPhyXMini0244.ChainMotionRegime}
  The dynamical regime described in the problem statement
\end{definition}
\begin{proof}
  This is the declared model definition of Chain Motion Regime.
\end{proof}
\begin{definition}[Force Term Treatment]
  \label{def:physics:phyx-mini-0244:phyxminiproblems-problemphyxmini0244-forcetermtreatment}
  \lean{PhyXMiniProblems.ProblemPhyXMini0244.ForceTermTreatment}
  Whether a force term is retained in, or idealized out of, the model
\end{definition}
\begin{proof}
  This is the declared model definition of Force Term Treatment.
\end{proof}
\begin{definition}[Rolling Chain Figure]
  \label{def:physics:phyx-mini-0244:phyxminiproblems-problemphyxmini0244-rollingchainfigure}
  \lean{PhyXMiniProblems.ProblemPhyXMini0244.RollingChainFigure}
  \uses{def:physics:phyx-mini-0244:phyxminiproblems-problemphyxmini0244-horizontaldirection, def:physics:phyx-mini-0244:phyxminiproblems-problemphyxmini0244-loopprofile, def:physics:phyx-mini-0244:phyxminiproblems-problemphyxmini0244-supportingsurfaceorientation, def:physics:phyx-mini-0244:phyxminiproblems-problemphyxmini0244-groundfeaturelabel}
  Qualitative information read from the supplied bitmap. The image gives no numerical radius, bump size, or separation, so none is invented here
\end{definition}
\begin{proof}
  This is the declared model definition of Rolling Chain Figure.
\end{proof}
\begin{definition}[Rolling Chain Setup]
  \label{def:physics:phyx-mini-0244:phyxminiproblems-problemphyxmini0244-rollingchainsetup}
  \lean{PhyXMiniProblems.ProblemPhyXMini0244.RollingChainSetup}
  \uses{def:physics:phyx-mini-0244:phyxminiproblems-problemphyxmini0244-chainradius, def:physics:phyx-mini-0244:phyxminiproblems-problemphyxmini0244-linearmassdensity, def:physics:phyx-mini-0244:phyxminiproblems-problemphyxmini0244-angularspeed, def:physics:phyx-mini-0244:phyxminiproblems-problemphyxmini0244-forcemagnitude, def:physics:phyx-mini-0244:phyxminiproblems-problemphyxmini0244-chainmotionregime, def:physics:phyx-mini-0244:phyxminiproblems-problemphyxmini0244-forcetermtreatment, def:physics:phyx-mini-0244:phyxminiproblems-problemphyxmini0244-rollingchainfigure}
  The physical quantities in the rolling-chain experiment. The single field uniformLinearMassDensity is the constant density μ of the uniform chain. The three speed fields have distinct roles: centerOfMassSpeed is the magnitude v₀ printed next to the red arrow; materialCirculationSpeed is a link's speed around the loop relative to its center; transverseWaveSpeed is the propagation speed of a transverse disturbance in the tensioned chain. No field defines tension from any of these quantities
\end{definition}
\begin{proof}
  This is the declared model definition of Rolling Chain Setup.
\end{proof}
\begin{definition}[Matches Problem Statement]
  \label{def:physics:phyx-mini-0244:phyxminiproblems-problemphyxmini0244-matchesproblemstatement}
  \lean{PhyXMiniProblems.ProblemPhyXMini0244.MatchesProblemStatement}
  \uses{def:physics:phyx-mini-0244:phyxminiproblems-problemphyxmini0244-rollingchainsetup}
  Prose assumptions, including the explicitly stated high-speed idealization
\end{definition}
\begin{proof}
  This is the declared model definition of Matches Problem Statement.
\end{proof}
\begin{definition}[Matches Supplied Figure]
  \label{def:physics:phyx-mini-0244:phyxminiproblems-problemphyxmini0244-matchessuppliedfigure}
  \lean{PhyXMiniProblems.ProblemPhyXMini0244.MatchesSuppliedFigure}
  \uses{def:physics:phyx-mini-0244:phyxminiproblems-problemphyxmini0244-rollingchainsetup}
  Qualitative geometry and labels visible in the primary image
\end{definition}
\begin{proof}
  This is the declared model definition of Matches Supplied Figure.
\end{proof}
\begin{definition}[Has Physical Rolling Parameters]
  \label{def:physics:phyx-mini-0244:phyxminiproblems-problemphyxmini0244-hasphysicalrollingparameters}
  \lean{PhyXMiniProblems.ProblemPhyXMini0244.HasPhysicalRollingParameters}
  \uses{def:physics:phyx-mini-0244:phyxminiproblems-problemphyxmini0244-radiusreadout, def:physics:phyx-mini-0244:phyxminiproblems-problemphyxmini0244-linearmassdensityreadout, def:physics:phyx-mini-0244:phyxminiproblems-problemphyxmini0244-speedreadout, def:physics:phyx-mini-0244:phyxminiproblems-problemphyxmini0244-angularspeedreadout, def:physics:phyx-mini-0244:phyxminiproblems-problemphyxmini0244-rollingchainsetup}
  Positivity and nondegeneracy conditions for the high-speed rolling regime
\end{definition}
\begin{proof}
  This is the declared model definition of Has Physical Rolling Parameters.
\end{proof}
\begin{definition}[Satisfies Rolling Chain Laws]
  \label{def:physics:phyx-mini-0244:phyxminiproblems-problemphyxmini0244-satisfiesrollingchainlaws}
  \lean{PhyXMiniProblems.ProblemPhyXMini0244.SatisfiesRollingChainLaws}
  \uses{def:physics:phyx-mini-0244:phyxminiproblems-problemphyxmini0244-radiusreadout, def:physics:phyx-mini-0244:phyxminiproblems-problemphyxmini0244-linearmassdensityreadout, def:physics:phyx-mini-0244:phyxminiproblems-problemphyxmini0244-speedreadout, def:physics:phyx-mini-0244:phyxminiproblems-problemphyxmini0244-angularspeedreadout, def:physics:phyx-mini-0244:phyxminiproblems-problemphyxmini0244-forcereadout, def:physics:phyx-mini-0244:phyxminiproblems-problemphyxmini0244-rollingchainsetup}
  The governing physics, stated in every coherent unit system: rolling without slipping gives v₀ = ω R; material links circulate around the center at speed u = ω R; persistence of the moving chain shape matches the transverse-wave speed to that material circulation speed; the standard transverse-wave law for a tensioned uniform chain is c² = T / μ. The last law is deliberately written using the independent wave-speed field; none of these premises states the requested formula involving v₀
\end{definition}
\begin{proof}
  This is the declared model definition of Satisfies Rolling Chain Laws.
\end{proof}
\begin{lemma}[transverse Wave Speed eq center Of Mass Speed]
  \label{lem:physics:phyx-mini-0244:phyxminiproblems-problemphyxmini0244-transversewavespeed-eq-centerofmassspeed}
  \lean{PhyXMiniProblems.ProblemPhyXMini0244.transverseWaveSpeed_eq_centerOfMassSpeed}
  \uses{def:physics:phyx-mini-0244:phyxminiproblems-problemphyxmini0244-speedreadout, def:physics:phyx-mini-0244:phyxminiproblems-problemphyxmini0244-rollingchainsetup, def:physics:phyx-mini-0244:phyxminiproblems-problemphyxmini0244-satisfiesrollingchainlaws}
  Rolling kinematics and stable-shape propagation identify the transverse-wave speed with the center-of-mass speed. This lemma derives a speed relation only; it does not assert the requested tension
\end{lemma}
\begin{proof}
  The conclusion follows from the governing relations and figure readouts named in the statement of transverse Wave Speed eq center Of Mass Speed.
\end{proof}
\begin{definition}[Answer Choice]
  \label{def:physics:phyx-mini-0244:phyxminiproblems-problemphyxmini0244-answerchoice}
  \lean{PhyXMiniProblems.ProblemPhyXMini0244.AnswerChoice}
  Labels of the four displayed answer choices
\end{definition}
\begin{proof}
  This is the declared model definition of Answer Choice.
\end{proof}
\begin{definition}[Displayed Speed Scaling]
  \label{def:physics:phyx-mini-0244:phyxminiproblems-problemphyxmini0244-displayedspeedscaling}
  \lean{PhyXMiniProblems.ProblemPhyXMini0244.DisplayedSpeedScaling}
  The speed dependence printed in each choice. The source uses a generic coefficient symbol α; its powers have different dimensions across choices, so the alternatives are retained as symbolic scaling metadata rather than as ill-typed force equalities
\end{definition}
\begin{proof}
  This is the declared model definition of Displayed Speed Scaling.
\end{proof}
\begin{definition}[displayed Speed Scaling]
  \label{def:physics:phyx-mini-0244:phyxminiproblems-problemphyxmini0244-displayedspeedscaling-value}
  \lean{PhyXMiniProblems.ProblemPhyXMini0244.displayedSpeedScaling}
  \uses{def:physics:phyx-mini-0244:phyxminiproblems-problemphyxmini0244-answerchoice, def:physics:phyx-mini-0244:phyxminiproblems-problemphyxmini0244-displayedspeedscaling}
  Speed scaling associated with each displayed answer label
\end{definition}
\begin{proof}
  This is the declared model definition of displayed Speed Scaling.
\end{proof}
\begin{definition}[recorded Answer Choice]
  \label{def:physics:phyx-mini-0244:phyxminiproblems-problemphyxmini0244-recordedanswerchoice}
  \lean{PhyXMiniProblems.ProblemPhyXMini0244.recordedAnswerChoice}
  \uses{def:physics:phyx-mini-0244:phyxminiproblems-problemphyxmini0244-answerchoice}
  Answer label recorded in the source dataset
\end{definition}
\begin{proof}
  This is the declared model definition of recorded Answer Choice.
\end{proof}
% --- Archon declaration topology end ---
% --- Archon physics formalization source end ---
